\documentclass[10pt,a4paper]{amsart}
\usepackage[margin=19mm]{geometry}
\usepackage{amsmath,amssymb,mathtools}
\usepackage[scr=rsfs,cal=euler]{mathalfa}
\usepackage{xfrac}
\usepackage[unicode,hidelinks,bookmarks=false]{hyperref}
\newcommand{\ie}{i.e.\ }
\newcommand{\cf}{cf.\ }
\newcommand{\resp}{respectively\ }
\newcommand{\Subsection}{\subsection}
\providecommand{\nobreakdash}{\nobreak\hbox{-}}
\setlength{\emergencystretch}{2em}
\multlinegap0pt
\usepackage{mathtools}
\usepackage{amssymb}
\newtheorem{lemma}{Lemma}
\newtheorem{proposition}{Proposition}
\newcommand{\irm}{\mathrm{i}}
\title{On the generating series of the degree sequence\date{\vspace{-2\baselineskip}}}
\author{Quang-Khai Nguyen}
\date{Source: arXiv:2510.06142v2 (CC BY 4.0)}
\begin{document}
\maketitle

\noindent\emph{Source and licence.} On the generating series of the degree sequence\date{\vspace{-2\baselineskip}}. Version arXiv:2510.06142v2 (CC BY 4.0), distributed by arXiv under the Creative Commons Attribution 4.0 International licence, http://creativecommons.org/licenses/by/4.0/, as recorded in the arXiv licence metadata for this exact version and shown on the abstract page https://arxiv.org/abs/2510.06142v2. Full licence notice: \texttt{licenses/arxiv-2510.06142-CC-BY-4.0.txt}.\par\smallskip

\noindent\textit{Editorial scope.} Counted unit: the article's proposition functions (source0.tex lines 335--339). The article's complete original proof of it is delivered with it (source0.tex lines 340--363, one \texttt{proof} environment of 24 lines). No mathematics is written, repaired or re-derived: the statement and the proof are the article's own lines, in the article's own notation, and no retained line is substituted or re-flowed.  Retained uncounted as the source-local context the proof needs: lines 328--330.  Nothing was omitted from any retained span, so the package carries no omission record.  No retained line contains a citation, so the wrapper carries no bibliography and no bibliography entry.  No retained line contains a figure environment, an image-inclusion command or drawing code, so no image is delivered and the package declares no asset dependency.  Editorial formatting only, in addition: an AMS \texttt{amsart} preamble that restates the article's own declarations and macros used by the retained lines, namely the environments \texttt{lemma}, \texttt{proposition} on separate counters, together with the standard packages those lines need.  The retained environments print in this document as Lemma 1, Proposition 1 in order of appearance, whereas the article prints the same items with the numbering of its own full document.  The complete native source of the article as distributed in the arXiv e-print for this exact version is bundled unchanged as \texttt{sources/186586/source0.tex}.
\begin{lemma}\label{piecewise linear functions are Lipschitz}
   Every piecewise linear function is Lipschitz.
\end{lemma}

\begin{proposition}\label{Fourier coefficients of continuous piecewise linear functions}
     Let $g\colon \mathbb R^2\rightarrow \mathbb R$ be a piecewise linear function. Let $\chi$ be some non-zero complex number. We denote by $\theta$ its argument. Then the power series $\sum_{n\geq0}g(\mathrm{Re}(\chi^n),\mathrm{Im}(\chi^n))z^n$ converges inside the disk $\mathbb D(0,|\chi|^{-1})$ and can be written as 
     $$\sum_{m\in\mathbb Z}\frac{a_m}{1-e^
    {\irm m\theta}|\chi|z}$$
     where $a_m\in\mathbb C$ such that $\sum_{m\in\mathbb Z}|a_m|<\infty$. Further, there is a linear recurrence sequence $\{d_m:m\in\mathbb Z\},$ such that $a_m=\frac{d_m}{m^2-1}$ for all $m\neq\pm1$. \end{proposition}

\begin{proof} 
  Since $g$ is piecewise linear, we have $|g(\mathrm{Re}(\chi^n),\mathrm{Im}(\chi^n))|= O(|\chi|^n)$ which yields that $\sum_{n\geq0}g(\mathrm{Re}(\chi^n),\mathrm{Im}(\chi^n))z^n$ converges inside the disk $\mathbb D(0,|\chi|^{-1})$. 
  
  It is convenient to identify $\mathbb R^2$ with $\mathbb C$ to make use of the conformal structure on $\mathbb C$, and view $g$ as a  1-homogeneous continuous function on $\mathbb C$. Then  $$\sum_{n\geq0}g(\mathrm{Re}(\chi^n),\mathrm{Im}(\chi^n))z^n=\sum_{n\geq0}|\chi|^ng(e^{\irm n\theta})z^n$$
   inside the disk $\mathbb D(0,|\chi|^{-1}).$

We decompose the function $g|_{\mathbb S^1}$ in the Hilbertian basis consisting of $\{e^{\irm mx}:m\in\mathbb Z\}$ in $L^2(\mathbb S^1)$. Since $g|_{\mathbb S^1}$ is a piecewise linear function on $\mathbb S^1$, it is Lipschitz by Lemma~\ref{piecewise linear functions are Lipschitz}. Therefore, its Fourier series converges uniformly to $g|_{\mathbb S^1}$. In addition, since $g|_{\mathbb S^1}$ is piecewise linear, $g|_{\mathbb S^1}$ belongs to the Sobolev space $W^{1,2}(\mathbb S^1)\coloneqq \{f\in L^2(\mathbb S^1):f'\in L^2(\mathbb S^1) \}$. Thus, for $x\in[0,2\pi]$, we may write $$g(e^{\irm  x})=\sum_{m\in\mathbb Z}a_{m}e^{\irm mx}$$
for $a_m\in\mathbb C$ and we have $\sum_{m\in\mathbb Z}m^2|a_m|^2<\infty$ by Parseval's identity. It follows that $\sum_{m\in\mathbb Z}|a_m|<\infty$. Therefore, for $z\in\mathbb{D}(0,|\chi|^{-1})$, we obtain
\begin{align*}
    \sum_{n\geq0}|\chi|^ng(e^{\irm n\theta})z^n&=\sum_{m\in\mathbb Z}a_m\sum_{n\geq0}|\chi|^ne^{\irm mn\theta}z^n=\sum_{m\in\mathbb Z}a_m\frac{1}{1-e^
    {\irm m\theta}|\chi|z}.
\end{align*} 

Now we compute $a_m$. Since $g$ is piecewise linear, we can write 
$g(z)=b_jz+c_j\overline{z}$ when $\arg z\in[\alpha_j,\alpha_{j+1})$, where $0=\alpha_0<\alpha_1<\ldots<\alpha_{t}=2\pi$ are the break points on the unit circle $\mathbb S^1$ for some $t\in\mathbb N$. For $m\neq\pm1$, we have
\begin{align*}
  2\pi a_{m}&=\int_{0}^{2\pi}g(e^{\irm x})e^{-\irm m x}dx=\sum_{j=0}^{t-1}\int_{\alpha_j}^{\alpha_{j+1}}b_je^{\irm  (-m+1)x}+c_je^{\irm  (-m-1)x}dx\\&=\sum_{j=0}^{t-1}b_j\frac{e^{\irm (-m+1)\alpha_{j+1}}-e^{\irm (-m+1)\alpha_{j}}}{\irm(-m+1)}+\sum_{j=0}^{t-1}c_j\frac{e^{\irm (-m-1)\alpha_{j+1}}-e^{\irm (-m-1)\alpha_{j}}}{\irm(-m-1)}.
\end{align*}
For each $j=0,\ldots,t$, we denote $A_j=e^{\irm \alpha_j}$, then\begin{align*}
    a_{m}&=\frac{\irm(m+1)\sum_{j=0}^{t-1}b_j(A_{j+1}^{-m+1}-A_j^{-m+1})+\irm(m-1)\sum_{j=0}^{t-1}c_j(A_{j+1}^{-m-1}-A_j^{-m-1})}{2\pi(m^2-1)}\\
    &=\frac{d_m}{m^2-1}
\end{align*}
where $d_m$ is a power sum in $m$; thus, the sequence $\{d_m:m\in\mathbb Z\}$ satisfies a linear recurrence relation as wanted.
\end{proof}

\end{document}
